\section{问题重述与研究目标}

\subsection{研究背景}

视频动作识别需要在连续图像中发现具有判别力的空间外观与时间运动模式。与单帧图像分类相比，
视频模型还必须处理帧数、运动跨度、背景变化、遮挡和长动作阶段缺失等问题。高容量的三维
卷积网络能够直接建模时空联合关系，但其激活和参数规模通常较大；二维骨干结合时序聚合可以
复用成熟的图像预训练权重，并通过共享空间编码器降低单帧计算成本。

本课题在 UCF101\cite{ucf101}的 101 类动作数据上，限定训练阶段 CUDA 显存峰值不超过 2 GiB，研究在有限
资源下应如何分配空间分辨率、时间长度和模型容量。课题不是简单地追求最高精度，而是要在同一
数据划分、输入预算和评估协议下比较不同时空建模路线，并给出可复核的效率结论。

\subsection{研究问题}

题目包含四项相互关联的研究内容。

第一，比较二维骨干加时序聚合与轻量三维卷积两条路线。二维路线采用 TSN 式分段采样\cite{tsn}，在共享
空间编码器上提取每一帧的特征，再用平均池化或时序注意力得到视频表示；三维路线采用朴素
R3D-18 作为基线\cite{i3d}，并采用带时空分解卷积的 R(2+1)D-18 作为轻量对照\cite{r21d}。

第二，研究均匀、稠密和运动自适应三种帧采样策略。三种策略固定使用同一个 TSN-2D 模型、
8 帧输入和 $112\times112$ 分辨率，只改变帧的位置选择方式，并记录额外解码和评分耗时。

第三，应用并评估混合精度、梯度检查点和梯度累积等显存优化技术。所有配置同时记录
\texttt{cuda max memory allocated} 与 \texttt{cuda max memory reserved} 两个峰值，
二者都不超过 2048 MiB 才视为满足预算。

第四，对类别代表性错误进行归因，包括 101 类混淆矩阵、分类别召回、低召回类别、主要混淆对
和动作语义组准确率，并据等成本的时间优先、均衡、空间优先实验回答“时间维还是空间维更值得
投入”。

\subsection{验收指标}

主指标为官方测试集 3783 个视频上的视频级 Top-1 准确率，辅助指标为 Top-5 准确率。候选路线
相对同预算朴素 R3D-18 的绝对提升应不少于 4 个百分点。每个核心头部配置使用 0、1、2 三个
训练种子，报告均值、样本标准差和同种子配对差值。若目标未达到，须保留负结果并分析原因，
不得通过减弱基线或更换测试集制造提升。
